\chapter[PROBLEM STATEMENT AND REQUIREMENTS]{Problem Statement and Requirements}

This chapter outlines the problem to be solved by TrustFin AI and converts it into application requirements. It outlines the user journeys, the functions available for each loan product and the technical limitations that influence the solution. It also separates the purpose of the decision support from the functions that would need a production lending environment.

\section{Problem Statement}
\subsection{Problem Definition}
A binary loan prediction returns a result without explaining how the model reached the result. This is not a strong basis for discussion with an applicant or a reviewing officer. A decision might need a closer look even if the classifier is doing well on the test data, due to the values submitted, the cost of the decision or the explanation behaviour.

The issue to be solved in this project is to make that individual assessment inspectable. The application should display the model output, include local reasons, compute defined checks on the local reasons and store the application context. It should also differentiate between prediction probability and explanation quality, so that one confident number does not take over the review.

\subsection{Limitations of a Prediction-Only Workflow}
Suppose there are two applicants who have been labeled with the same approval label. You can be able to have a comfortable repayment ratio and have a good credit score. The other one might be near to the model boundary or have monthly commitments that are considered too high by the application. These differences are not apparent when the common label is displayed.

This is also true of explanations. The feature names can be the same, but the directions can be different in two methods. A chart does not display whether the reasons chosen are stable in the presence of a given change. TrustFin AI solves these review requirements by providing the decision, financial summary, explanations and diagnostics as separate components of one response.

\section{Users and Functional Requirements}
\subsection{Users and Review Needs}
The applicant process starts with picking a product and filling out a form. The outcome should convey the decision, the key considerations and the repayment figures without assuming an understanding of the model internals. Here the submitted-form view is helpful: the user can see if the assessment has taken the information they were looking to provide.

The bank-officer path should be fleshed out. An officer can view SHAP and LIME contributions, compare the leading factors, view the four diagnostics and open a previously saved evaluation. Batch upload allows you to view multiple applications without having to fill out the same form multiple times. These trips share the same prediction service, but have varying levels of detail in the interface.

The role selection is a navigation feature. History endpoints are not currently restricted by account and the selected role is not authenticated. The project thus relies on local or controlled demonstration data and considers secure multi-user operation as a separate deployment requirement.

\subsection{Functional Requirements}
The principal requirements are related to the current application in Table~\ref{tab:functional-requirements}. Each requirement is a user visible function or a step needed to create it.

\begin{table}[htbp]\centering\small
\caption{Functional requirements}\label{tab:functional-requirements}
\begin{tabularx}{\textwidth}{p{1.3cm}X}\toprule
ID & Required function\\\midrule
FR-01 & Select a home-loan or personal-loan application flow.\\
FR-02 & Normalise application values and apply the product's financial feature calculations.\\
FR-03 & Return an approval or rejection prediction and model probability.\\
FR-04 & Generate local SHAP and LIME feature contributions.\\
FR-05 & Calculate the four implemented diagnostics and weighted Trust Score.\\
FR-06 & Present applicant and bank-officer result views over the shared assessment.\\
FR-07 & Save and retrieve evaluations, including their submitted inputs and explanation data.\\
FR-08 & Process CSV or JSON batches and identify unusable rows with reasons.\\
FR-09 & Generate an application PDF for preview or download.\\
FR-10 & Recalculate instalment and affordability values in the what-if calculator.\\\bottomrule
\end{tabularx}\end{table}

The result of the one evaluation should be sufficient for the subsequent evaluation to use the same assessment. This includes the dossier reference, product, prediction, probabilities, explanations, component scores and solvency summary. Otherwise, if only part of the result is recomputed in another screen, it would result in conflicting accounts of the same application.

\subsection{Product-Specific Requirements}
A home loan is a loan that is used to purchase a home. The application gathers property value and area, accepts co-applicant income and calculates loan-to-value ratio. A personal loan does not have any collateral fields, but rather relies on the loan purpose and the credit card utilisation. It bases its FOIR on the income of the applicants.

The model's feature representation is enforced by separating product preprocessors. The model selected should be passed fields in the same order and size as it was trained with. Any change to a form label should not change the feature passed to the classifier without notifying the user, and the addition of a display field should not be confused with a trained predictor.

If the product is used for batch input, then the product specifies the expected row structure. A helpful answer will indicate which rows were used and which were not. The system keeps the row number and error reason so that the user can fix the source file rather than trying to guess the reason for the loss of a record.

\section{Non-Functional Requirements}
The main design requirement is separation of responsibilities. The frontend displays application data and results. The API orchestrates the request, the model registry is responsible for product pipelines, and explanations, scoring, persistence and PDF generation are done with separate modules. This type of structure makes it easier to review changes in one area.

Repeatability is also a requirement. LIME sampling uses a seed derived from the application, and stability uses a deterministic set of nearby financial inputs. The repeated evaluation should be repeatable in the same model and software environment. This is not a library version record or proof of the accuracy of the methods, but it does make it easier to investigate the same input.

There should also be a functional local configuration for the application. SQLite offers persistence without a hosted database, and PostgreSQL can be configured via the environment. The typed request and response schemas ensure consistency of communication between the client and the server. There are no concurrency limits, response-time guarantees or service availability benchmarks.

\section{Software and Hardware Environment}
The backend needs Python, FastAPI, Pydantic, SQLAlchemy, pandas, NumPy, scikit-learn, XGBoost, SHAP and LIME. Frontend is built with React, Vite, Axios, Tailwind CSS and chart components. A browser, a Python process and the Node.js development toolchain provide the basic development setup.

The prerequisites for the project are Python 3.10 or higher and Node.js 18 or higher. The package manifests are the definitive dependency information, and the actual installed packages must meet the requirements. Most Python dependencies do not have strict version constraints, but rather a lower bound.

There is no tested hardware capacity or simultaneous user limit specified. The development workload consists of small tabular models, explanation sampling and local data storage. Measured latency, memory usage and concurrent-request tests on the target machine would be needed for a production capacity estimate.

\section{Acceptance and Operating Boundaries}
An application path is finished if the chosen product is able to receive a usable record, return a prediction and explanations, show the diagnostic values, and provide the history or export action. The verification chapter outlines how these functions can be verified without mixing up software behaviour with predictive accuracy.

The main limitations are the use of synthetic data, one saved train--test split per product and custom explanation-score definitions, as well as no authenticated access control. The financial ratios are not a universal lending policy, but rather project assumptions. The output from the model and the solvency summary displayed in the current route are still separate.

\par\medskip\noindent
In this chapter, the review problem has been identified, two user journeys have been defined and the functions that each product must perform have been identified. It has also established the software environment and operating limits that the design must adhere to.
